% ============================================================
% LEGENDAS (captions) de figuras, tabelas e outros flutuantes
% ============================================================

% Configura o estilo global das legendas
\captionsetup{
    font={footnotesize},   % 10pt em documento 12pt
    labelfont={bf},        % Rótulo (ex: "Figura 1") em negrito
    skip=0.3cm,            % Espaço entre o objeto e a legenda
    justification=centering, % Legenda centralizada
    singlelinecheck=true   % Centraliza automaticamente quando cabe em uma linha
}
\captionsetup[table]{skip=0.3cm}      % Espaço da legenda em tabelas
\captionsetup[figure]{skip=0.3cm}     % Espaço da legenda em figuras
\captionsetup[longtable]{skip=0.3cm}  % Espaço da legenda em tabelas longas

% ============================================================
% COMANDO CUSTOMIZADO: \fontetabela{texto}
% Exibe a fonte/referência abaixo de tabelas e figuras
% ============================================================
\providecommand{\fontetabela}[1]{%
    \par\vspace{0.3cm}%                          % Espaço acima do texto de fonte
    {\centering\fontsize{10pt}{10pt}\selectfont #1\par}% % Texto centralizado em 10pt
}

% ============================================================
% FONTE PRINCIPAL: Times New Roman (fonte do sistema)
% ============================================================
\setmainfont{Times New Roman} % Usa a fonte Times New Roman instalada no sistema operacional

% ============================================================
% MARGENS — padrão ABNT (NBR 14724)
% ============================================================
\geometry{
    a4paper,    % Tamanho do papel A4
    top=3cm,    % Margem superior
    left=3cm,   % Margem esquerda (encadernação)
    right=2cm,  % Margem direita
    bottom=2cm  % Margem inferior
}

\setstretch{1.5}               % Espaçamento entre linhas 1,5 (padrão ABNT)
\setlength{\parindent}{1.25cm} % Recuo de parágrafo 1,25 cm (padrão ABNT)

% ============================================================
% FORMATAÇÃO DOS TÍTULOS DE CAPÍTULOS E SEÇÕES
% ============================================================

% Capítulo: 14pt, negrito, maiúsculas, numerado
\titleformat{\chapter}[hang]
    {\normalfont\fontsize{14}{16}\bfseries\uppercase} % Estilo do título
    {\thechapter}  % Número do capítulo
    {1em}          % Espaço entre número e texto
    {}
% Espaçamentos: {recuo esquerdo}{espaço antes}{espaço depois}
\titlespacing*{\chapter}{0pt}{12pt}{12pt}

% Seção: 12pt, negrito, numerada
\titleformat{\section}[hang]
    {\normalfont\fontsize{12}{12}\bfseries}
    {\thesection}  % Número da seção
    {1em}
    {}
\titlespacing*{\section}{0pt}{24pt}{6pt}

% Subseção: 12pt, negrito, numerada
\titleformat{\subsection}[hang]
    {\normalfont\fontsize{12}{12}\bfseries}
    {\thesubsection} % Número da subseção
    {1em}
    {}
\titlespacing*{\subsection}{0pt}{18pt}{4pt}

% ============================================================
% RENOMEAR RÓTULOS para português (padrão ABNT)
% ============================================================
% Troca "Table"/"List of Tables" por "Quadro"/"Lista de Quadros"
\addto\captionsbrazil{%
    \renewcommand{\tablename}{Quadro}%
    \renewcommand{\listtablename}{Lista de Quadros}%
}

% ============================================================
% SUMÁRIO — formatação manual conforme ABNT
% ============================================================
\renewcommand{\contentsname}{}  % Remove o título automático "Sumário" (inserido manualmente)
\setlength{\cftbeforetoctitleskip}{0pt}
\setlength{\cftaftertoctitleskip}{0pt}
\renewcommand{\cfttoctitlefont}{\normalfont\fontsize{0.1pt}{0pt}\selectfont} % Título sem altura: fonte 0.1pt, entrelinha 0pt

% Fonte das entradas de capítulo no sumário (negrito, conforme imagem)
\renewcommand{\cftchapfont}{\normalfont\fontsize{12}{12}\bfseries\selectfont}
\renewcommand{\cftchappagefont}{\normalfont\fontsize{12}{12}\bfseries\selectfont}   % Número de página do capítulo

% Fonte das entradas de seção no sumário
\renewcommand{\cftsecfont}{\normalfont\fontsize{12}{12}\selectfont}
\renewcommand{\cftsecpagefont}{\normalfont\fontsize{12}{12}\selectfont}    % Número de página da seção

% Fonte das entradas de subseção no sumário
\renewcommand{\cftsubsecfont}{\normalfont\fontsize{12}{12}\selectfont}
\renewcommand{\cftsubsecpagefont}{\normalfont\fontsize{12}{12}\selectfont} % Número de página da subseção

% Pontinhos de condução (leaders) — estilo da imagem
\renewcommand{\cftchapdotsep}{\cftdotsep}   % Ativa pontinhos para capítulos
\renewcommand{\cftchapleader}{\cftdotfill{\cftchapdotsep}}
\renewcommand{\cftsecleader}{\cftdotfill{\cftdotsep}}
\renewcommand{\cftsubsecleader}{\cftdotfill{\cftdotsep}}

% Espaçamento vertical entre entradas do sumário (compacto)
\setlength{\cftbeforechapskip}{0pt}    % Antes de cada capítulo
\setlength{\cftbeforesecskip}{0pt}     % Antes de cada seção
\setlength{\cftbeforesubsecskip}{0pt}  % Antes de cada subseção

% Remove indentação de seções e subseções no sumário
\setlength{\cftsecindent}{0pt}
\setlength{\cftsubsecindent}{0pt}
\setlength{\cftsubsubsecindent}{0pt}

% ============================================================
% LISTA DE FIGURAS e LISTA DE QUADROS — formatação manual
% ============================================================
\renewcommand{\listfigurename}{}  % Remove título automático da lista de figuras
\renewcommand{\listtablename}{}   % Remove título automático da lista de quadros

% Configurações da lista de figuras (lof = list of figures)
\setlength{\cftbeforeloftitleskip}{0pt}
\setlength{\cftafterloftitleskip}{0pt}
\renewcommand{\cftloftitlefont}{\normalfont\fontsize{0.1pt}{0pt}\selectfont} % Título sem altura

% Configurações da lista de tabelas/quadros (lot = list of tables)
\setlength{\cftbeforelottitleskip}{0pt}
\setlength{\cftafterlottitleskip}{0pt}
\renewcommand{\cftlottitlefont}{\normalfont\fontsize{0.1pt}{0pt}\selectfont} % Título sem altura

% Pontinhos de condução nas listas de figuras e quadros
\renewcommand{\cftfigleader}{\cftdotfill{\cftdotsep}}   % Lista de figuras: com pontinhos
\renewcommand{\cfttableader}{\cftdotfill{\cftdotsep}}   % Lista de quadros: com pontinhos

% Fonte das entradas nas listas de figuras e quadros
\renewcommand{\cftfigfont}{\normalfont\fontsize{12}{12}\selectfont}
\renewcommand{\cftfigpagefont}{\normalfont\fontsize{12}{12}\selectfont}  % Número de página
\renewcommand{\cfttabfont}{\normalfont\fontsize{12}{12}\selectfont}
\renewcommand{\cfttabpagefont}{\normalfont\fontsize{12}{12}\selectfont}  % Número de página

% Prefixo e sufixo do número nas listas (ex: "Figura 1 –")
\renewcommand{\cftfigpresnum}{Figura }   % Texto antes do número na lista de figuras
\setlength{\cftfignumwidth}{5.5em}       % Largura reservada para "Figura 1 --"

\renewcommand{\cfttabpresnum}{Quadro }   % Texto antes do número na lista de quadros
\setlength{\cfttabnumwidth}{5.5em}       % Largura reservada para "Quadro 1 --"

% Espaçamento vertical entre entradas nas listas (compacto, igual ao sumário)
\setlength{\cftbeforefigskip}{0pt}  % Antes de cada figura
\setlength{\cftbeforetabskip}{0pt}  % Antes de cada quadro

% ============================================================
% COMANDO CUSTOMIZADO: \titulodelista{texto}
% Exibe o título centralizado e em negrito acima de cada lista
% ============================================================
\providecommand{\titulodelista}[1]{%
    \begin{center}%
        \fontsize{14pt}{14pt}\selectfont\bfseries #1% % 14pt, negrito, centralizado
    \end{center}%
    \vspace{4pt}% % Espaço após o título
}

% ============================================================
% PAGINAÇÃO — número da página no canto superior direito
% ============================================================
\pagestyle{fancy}          % Ativa o estilo customizado do pacote fancyhdr
\fancyhf{}                 % Limpa todos os campos de cabeçalho e rodapé
\fancyhead[R]{\thepage}    % Número da página no canto direito do cabeçalho
\renewcommand{\headrulewidth}{0pt} % Remove a linha horizontal do cabeçalho

% Reaplica o mesmo estilo nas páginas "plain" (ex: primeira página de capítulo)
\fancypagestyle{plain}{
    \fancyhf{}
    \fancyhead[R]{\thepage}
    \renewcommand{\headrulewidth}{0pt}
}

% Remove a dependência do contador de figuras/tabelas em relação ao capítulo
% (numeração contínua: Figura 1, 2, 3... em vez de Figura 1.1, 1.2...)
\counterwithout{figure}{chapter}
\counterwithout{table}{chapter}

% ============================================================
% FLOAT PERSONALIZADO: tabela (distinto de quadro)
% ============================================================
\DeclareFloatingEnvironment[
    name=Tabela,
    listname={Lista de Tabelas},
    within=none
]{tabela}

\renewcommand{\cfttabelapresnum}{Tabela }
\setlength{\cfttabelanumwidth}{5.5em}
\renewcommand{\cfttabelafont}{\normalfont\fontsize{12}{12}\selectfont}
\renewcommand{\cfttabelapagefont}{\normalfont\fontsize{12}{12}\selectfont}
\renewcommand{\cfttabelaleader}{\cftdotfill{\cftdotsep}}
\setlength{\cftbeforetabelaskip}{0pt}
